% Section 7.3: case (ii), equal scan cost. Facts: log items 88-89, 94-95, 103, 108; figures F8c/F8d, F10a.
\subsection{Case (ii): the same scan budget}\label{sec:case-ii}
\begin{notes}
\item Pre-registered Fisher test (all five found, adaptive batches vs static): 2/5 vs 1/8, 1/5 vs 0/7, 0/4 vs 1/7 -- not significant.
\item Diagnosis: D4 is located correctly by the adaptive runs but too weak (1.4--2.5\,mm); the round-2 designs chase false thin spots of the 20-point model.
\item Re-inversion (adaptive points, one pass) vs static at 1\,\% rate (exploratory): 4.20 / 4.00 / 4.00 vs 3.88 / 3.43 / 3.43, all five 4/14 vs 2/22, MWU $p=0.022$; paired with the same start 5/4/1 ($p=0.11$); fixed-amplitude rule: no gain. False indications lower than as run in 12 of 14 re-inversions.
\item Second criterion ($-\overline{\log I}$): no gain anywhere (report, do not hide).
\item \note{24 Nov}: simulation design (designed beforehand) vs static and vs re-inverted adaptive: if it is as good, the gain is from design, not adaptivity. Decide whether to keep.
\item \textbf{24 Nov results}: 30 points designed beforehand (3 per shot): 1\,\% 4.67 / 3.33 / 2.67, v3 4.67 / 3.33 / 3.00 vs static 3.88 / 3.43 / 3.43: no gain except with 4 shots (same start 3/2/4). Re-inverted adaptive vs designed beforehand, same start: 4/4/1 (1\,\%), 5/3/1 (5\,\%, $p=0.11$), 2/4/3 (v3): a small point for designing with the data.
\item Second criterion: move to the appendix (no gain anywhere).
\end{notes}
